\documentclass[pdflatex,sn-nature]{sn-jnl}% Style for submissions to Nature Portfolio journals

%%%% Standard Packages
\usepackage{graphicx}%
\usepackage{multirow}%
\usepackage{amsmath,amssymb,amsfonts}%
\usepackage{amsthm}%
\usepackage{mathrsfs}%
\usepackage[title]{appendix}%
\usepackage{xcolor}%
\usepackage{textcomp}%
\usepackage{manyfoot}%
\usepackage{booktabs}%
\usepackage{algorithm}%
\usepackage{algorithmicx}%
\usepackage{algpseudocode}%
\usepackage{listings}%
\theoremstyle{thmstyleone}%
\newtheorem{theorem}{Theorem}%  meant for continuous numbers
%%\newtheorem{theorem}{Theorem}[section]% meant for sectionwise numbers
%% optional argument [theorem] produces theorem numbering sequence instead of independent numbers for Proposition
\newtheorem{proposition}[theorem]{Proposition}%
%%\newtheorem{proposition}{Proposition}% to get separate numbers for theorem and proposition etc.

\theoremstyle{thmstyletwo}%
\newtheorem{example}{Example}%
\newtheorem{remark}{Remark}%

\theoremstyle{thmstylethree}%
\newtheorem{definition}{Definition}%

\raggedbottom
%%\unnumbered% uncomment this for unnumbered level heads

\begin{document}

\title[Modular and Scalable Analysis of Proteomics Data with alphapepttools]{Modular and Scalable Analysis of Proteomics Data with alphapepttools}

%%=============================================================%%
%% GivenName  -> \fnm{Joergen W.}
%% Particle  -> \spfx{van der} -> surname prefix
%% FamilyName  -> \sur{Ploeg}
%% Suffix  -> \sfx{IV}
%% \author*[1,2]{\fnm{Joergen W.} \spfx{van der} \sur{Ploeg}
%%  \sfx{IV}}\email{iauthor@gmail.com}
%%=============================================================%%

\author[1,2]{\fnm{Vincenth} \sur{Brennsteiner}}

\author[1,2]{\fnm{Lucas} \sur{Diedrich}}

\author[1]{\fnm{Shani} \sur{Ben-Moshe}}

\author[1]{\fnm{Magnus} \sur{Schwörer}}

\author[2]{\fnm{scverse proteomics interest group}}

\author*[1]{\fnm{Matthias} \sur{Mann}}\email{mmann@biochem.mpg.de}


\affil[1]{\orgdiv{Department for Proteomics and Signal Transduction}, \orgname{Max Planck Institute of Biochemistry}, \orgaddress{\street{Am Klopferspitz 18}, \city{Planegg}, \postcode{821252}, \state{Bavaria}, \country{Germany}}}
\affil[2]{\orgdiv{Proteomics Interest Group}, \orgname{scverse}}

\abstract{Abstract.}

\keywords{LC/MS Proteomics, Data Analysis}

%%\pacs[JEL Classification]{D8, H51}

%%\pacs[MSC Classification]{35A01, 65L10, 65L12, 65L20, 65L70}

\maketitle

\section{Introduction}\label{sec:introduction}
% 500 words

\section{Results}\label{sec:results}
% \section, \subsection, \subsubsection
\subsection{Application of alphapepttools across diverse proteomics workflows}
To illustrate practical use cases, we reanalyzed selected datasets to demonstrate the versatility of alphapepttools across diverse data types, including low-input single-cell proteomics, matched proteomics and transcriptomics, and peptidomics analyses.
For low-input single-cell proteomics, we revisited a dataset from Weiss et al. (2026) [ https://www.nature.com/articles/s42255-026-01459-2] that profiled the proteomes of individual human hepatocytes along the porto-central anatomical axis using Deep Visual Proteomics (DVP) [https://www.nature.com/articles/s41587-022-01302-5]. Using alphapepttools' built-in functionality, we reproduced the core elements of the published analysis workflow, including quality control visualization, data filtering, and dimensionality reduction via PCA (Fig. XXX).
Notably, we extended the downstream analysis by inferring transcription factor activities associated with the different spatial layers. To this end, Decoupler package [https://academic.oup.com/bioinformaticsadvances/article/2/1/vbac016/6544613?login=true] from the scverse ecosystem [DOI: 10.1038/s41587-023-01733-8 ]. The AnnData-based enabled a streamlined interoperability with widely adopted computational tools.
Taken together, alphapepttools' built-in functions and its inherent compatibility with available analysis tools facilitates efficient and biologically meaningful interpretation of proteomics datasets.

\section{Discussion}\label{sec:discussion}


\backmatter

\bmhead{Supplementary information}

If your article has accompanying supplementary file/s please state so here.


\bmhead{Acknowledgements}
% Acknowledgements are not compulsory. Where included they should be brief. Grant or contribution numbers may be acknowledged.
% Please refer to Journal-level guidance for any specific requirements.

\section*{Declarations}
\subsection*{Funding}
Not applicable

\subsection*{Competing interests}

\subsection*{Ethics approval and consent to participate}
Not applicable

\subsection*{Consent for publication}
% All authors read the manuscript and agreed to its publication.

\subsection*{Data availability}

\subsection*{Materials availability}
Not applicable

\subsection*{Code availability}

\subsection*{Author contribution}


\begin{appendices}

  \section{Extended Data}\label{secA1}

  %   An appendix contains supplementary information that is not an essential part of the text itself but which may be helpful in providing a more comprehensive understanding of the research problem or it is information that is too cumbersome to be included in the body of the paper.

\end{appendices}

%%===========================================================================================%%
%% If you are submitting to one of the Nature Portfolio journals, using the eJP submission   %%
%% system, please include the references within the manuscript file itself. You may do this  %%
%% by copying the reference list from your .bbl file, paste it into the main manuscript .tex %%
%% file, and delete the associated \verb+\bibliography+ commands.                            %%
%%===========================================================================================%%

\bibliography{bibliography}% common bib file
%% if required, the content of .bbl file can be included here once bbl is generated
%%\input sn-article.bbl

\end{document}
